Co-authorship is not a technique. It is what an interaction looks like 
when both people are genuinely present and neither is resolving meaning 
before it has formed between them.

It is recognisable not by what either person does, but by what the 
interaction produces. Meaning does not arrive from one side and land 
on the other. It accumulates in the space between --- built by both, 
owned by neither, carrying the fingerprints of each.

This is what distinguishes it from assertion with flipped dynamics. 
When only one person leaves meaning open and the other simply fills it, 
the structure has not changed --- only the direction. One person is still 
operating on the other. Co-authorship requires that both people are 
simultaneously incomplete, simultaneously reaching. Neither is the 
architect. Neither is the audience.

What this looks like varies. It can look like a silence neither person 
moves to fill, because neither needs to. It can look like a thought 
completed by the wrong person --- the one who didn't start it finishing 
it more accurately than the one who did. It can look like physical 
presence that neither initiated and neither interrupted. It can look 
like two people arriving at the same unspoken conclusion without 
either having directed the other there.

The common thread is not a behaviour. It is an absence: the absence 
of either person managing the interaction toward an outcome. When that 
absence is genuine on both sides, meaning builds on its own. 

Co-authorship is the symptom. Generative Incompletion, operating 
from both sides, is the condition that produces it.
\subsection{The Completion Bond}
When meaning is built mutually, the investment it produces is different 
in kind from anything delivered complete. A meaning that arrived from 
outside can be held at a distance, examined, agreed with or rejected. 
A meaning that was built from inside cannot. It carries the builder's 
own material. They cannot be indifferent to what they made.

This is the Completion Bond: not an attachment to the other person, 
but an attachment to the meaning itself --- and through it, to the 
interaction that made the meaning possible. It does not require 
intensity to form. It requires only that the meaning was genuinely 
incomplete before they moved toward it, and that they moved of their own 
volition.

What distinguishes it from ordinary interest is its durability. 
Interest produced by delivered meaning closes when the meaning is 
fully received. The Completion Bond does not close --- because the 
meaning was never fully delivered from outside. They wrote part of it. 
That part stays with them, because it came from them.

This also explains the reciprocal disclosure it tends to produce. 
When meaning has been built mutually, the psychological position 
of both people shifts. Neither was operating on the other. Neither 
was the target of a move. Something real was built in shared space, 
and the nervous system of each person reads the other as safe --- 
not because trust was established through disclosure, but because 
the interaction itself was not coercive. Her guard does not lower 
as a decision. It lowers as a consequence.

\begin{calloutbox}
Delivered meaning produces a receiver.\\
Built meaning produces a participant.\\[0.4em]
The receiver can walk away from what they were given.\\
The participant carries what they helped make.
\end{calloutbox}
\subsection{Let Her Move Toward You}

Co-authorship cannot be forced. The man who presses for completion --- 
who crowds the space before the receiver has had time to fill it --- does not accelerate 
the mechanism. He replaces it. What arrives under pressure is compliance, 
not investment. They moved because the space closed, not because they chose to enter it.

The Completion Bond forms only when they move on their own terms. 
What they reache voluntarily is theirs. What they are directed towards is not. 
The distinction is not subtle --- it determines whether the meaning they 
carry after the interaction belongs to them or to the interaction itself.

This is why restraint is not a stylistic preference in this framework. 
It is a structural requirement. The space in which they completes the meaning 
must remain genuinely open --- not performed open, not strategically open, 
but open because the speaker did not reach past the natural edge of what had formed.

\begin{calloutbox}
Assertion delivers meaning complete.\\
Co-authorship leaves meaning open at the edge where it genuinely sits.\\[0.4em]
One produces a receiver. The other produces a participant.\\
Participation cannot be forced. It can only be made possible.
\end{calloutbox}

\subsection{The Collapse Condition}

The Completion Bond does not only fail to form. It collapses --- 
quietly, without announcement --- when the conditions that produced 
it are reversed.

The reversal does not require hostility or error. It requires only 
that one person begins resolving meaning before the other has finished 
building it. The shared space closes. What was mutual becomes directed. 
The interaction has not ended, but co-authorship has.

This happens in recognisable ways.

When meaning that has just landed is explained --- elaborated, softened, 
clarified --- the space in which they was still completing it is filled 
from outside. They was mid-construction. The construction is finished 
for them. What they were reaching towards, building, becomes something they were handed. 
The reaching stops.

When what they completed is named --- reflected back, confirmed, 
labelled --- the meaning moves from the shared space into explicit 
territory. It becomes an object between two people rather than 
something alive in the space they are both inside. It cannot be 
returned to what it was. The intimacy of built meaning depends on 
it remaining partially unspoken, even after it has been felt by both.

When one person begins moving before the other has finished --- 
pressing the next fragment, closing the distance, advancing the 
interaction --- the other person's motion converts from voluntary 
to reactive. They were choosing to move. They are now responding to 
pressure. These feel similar from the outside. They are not the 
same from the inside. The Completion Bond requires that both people 
remain genuinely free to stop. The moment that ceases to be true, 
what is forming is no longer co-authorship.

\begin{calloutbox}
The bond does not break dramatically.\\
It stops being real.\\[0.4em]
Explain after it lands --- they stops building.\\
Name what was built --- it becomes an object.\\
Move before the receiver finishes --- they stop choosing.\\[0.4em]
Co-authorship requires that both people\\
remain free to stop at any moment.\\
That freedom is not a courtesy.\\
It is the condition.
\end{calloutbox}